% Generado por sistema/backend ReporteTex.java -- no editar a mano.
% Datos: dataset de prueba (respuestas sinteticas)
\subsubsection*{C4. Interés en posgrado $\times$ oferta de posgrado FCyT}

\textit{Propósito.} Determinar si la percepción de que la oferta de posgrado de la FCyT es afín a la carrera depende del interés en realizar un posgrado. Respalda el criterio ARCU-SUR 2.4.1 / 1.1.5. Se aplica sobre: todos los titulados, con nivel de significancia $\alpha = 0.05$.

\textit{Paso 1. Contar con la muestra: tabla de contingencia} (frecuencia observada $FO_{ij}$, $N = 100$).

\begin{mitabla}{C4: frecuencia observada $FO_{ij}$ \datosprueba}
\begin{tabular}{lcc|c}
\toprule
Interés en posgrado & En desacuerdo & De acuerdo & Total fila \\
\midrule
Sí & 15 & 49 & 64 \\
No & 11 & 25 & 36 \\
\midrule
Total columna & 26 & 74 & $N=100$ \\
\bottomrule
\end{tabular}
\end{mitabla}

\textit{Paso 2. Calcular la frecuencia esperada} de cada celda, $FE_{ij} = (\text{total fila } i)(\text{total columna } j)/N$. Por ejemplo, para la celda ``Sí -- En desacuerdo'': $FE_{11} = 64 \times 26 / 100 = 16.6400$.

\begin{mitabla}{C4: frecuencia esperada $FE_{ij}$ \datosprueba}
\begin{tabular}{lcc}
\toprule
Interés en posgrado & En desacuerdo & De acuerdo \\
\midrule
Sí & $64\cdot26/100 = 16.64$ & $64\cdot74/100 = 47.36$ \\
No & $36\cdot26/100 = 9.36$ & $36\cdot74/100 = 26.64$ \\
\bottomrule
\end{tabular}
\end{mitabla}

\textit{Paso 3. Calcular el estadístico chi-cuadrado,} sumando $(FO_{ij}-FE_{ij})^2/FE_{ij}$ de cada celda.

\begin{mitabla}{C4: cálculo de $X_0^2$ celda por celda \datosprueba}
\begin{tabular}{L{5.4cm}cccc}
\toprule
Celda & $FO_{ij}$ & $FE_{ij}$ & $(FO_{ij}-FE_{ij})^2$ & $(FO_{ij}-FE_{ij})^2/FE_{ij}$ \\
\midrule
Sí -- En desacuerdo & 15 & 16.6400 & 2.6896 & 0.1616 \\
Sí -- De acuerdo & 49 & 47.3600 & 2.6896 & 0.0568 \\
No -- En desacuerdo & 11 & 9.3600 & 2.6896 & 0.2874 \\
No -- De acuerdo & 25 & 26.6400 & 2.6896 & 0.1010 \\
\midrule
\multicolumn{4}{r}{$X_0^2 = \sum$} & 0.6067 \\
\bottomrule
\end{tabular}
\end{mitabla}

\textit{Paso 4. Determinar el estadístico teórico.} Con $k = 2$ filas y $m = 2$ columnas, los grados de libertad son $(k-1)(m-1) = (2-1)(2-1) = 1$. De la tabla de chi-cuadrada: $X_{0.05,\,1}^2 = 3.841$.

\textit{Paso 5. Establecer la relación entre estadísticos:} $X_0^2 < X_{0.05,\,1}^2 \;\Rightarrow\; 0.6067 < 3.841$.

\textit{Paso 6. Conclusión.}

\begin{center}
\begin{tabular}{l}
Si ($X_0^2 < X_{0.05,\,1}^2$)\\
\quad \fbox{Independencia}\\
Sino\\
\quad Relación entre las variables\\
FinSi
\end{tabular}
\end{center}

No se puede rechazar que interés en posgrado y oferta fcyt afín sean independientes. \datosprueba

